\textbf{Texto de apoio para as questões 01 e 02}\par\medskip
\textbf{Velha Infância}\\

Você é assim
Um sonho pra mim
E quando eu não te vejo
Eu penso em você
Desde o amanhecer
Até quando eu me deito

Eu gosto de você
E gosto de ficar com você
Meu riso é tão feliz contigo
O meu melhor amigo é o meu amor

E a gente canta
E a gente dança
E a gente não se cansa
De ser criança
A gente brinca
Na nossa velha infância

Seus olhos, meu clarão
Me guiam dentro da escuridão
Seus pés me abrem o caminho
Eu sigo e nunca me sinto só

Você é assim
Um sonho pra mim
Quero te encher de beijos
Eu penso em você
Desde o amanhecer
Até quando eu me deito

Eu gosto de você
E gosto de ficar com você
Meu riso é tão feliz contigo
O meu melhor amigo é o meu amor

E a gente canta
A gente dança
A gente não se cansa
De ser criança
A gente brinca
Na nossa velha infância

Seus olhos, meu clarão
Me guiam dentro da escuridão
Seus pés me abrem o caminho
Eu sigo e nunca me sinto só

Você é assim
Um sonho pra mim
Você é assim

Você é assim
Um sonho pra mim
Você é assim

Você é assim
Um sonho pra mim
Você é assim

Você é assim
Um sonho pra mim
E quando eu não te vejo
Eu penso em você
Desde o amanhecer
Até quando eu me deito

Eu gosto de você
E gosto de ficar com você
Meu riso é tão feliz contigo
O meu melhor amigo é o meu amor

(Tribalistas)

O texto apresenta uma declaração de afeto e a alegria de compartilhar o cotidiano com alguém. Em um breve trecho, a voz da canção situa esse sentimento ao longo do dia:
\begin{center}\textit{Desde o amanhecer / Até quando eu me deito}\end{center}
Disponível em: \url{http://www.letras.mus.br/tribalistas/64146}. Acesso em: 11 ago. 2019.

\textbf{COLUNI 2020 -- Questão 01}\par\medskip Sobre o texto \textit{Velha Infância}, assinale a afirmativa CORRETA:\begin{enumerate}
\item O amor é tema abordado de forma simples e inerente ao ser humano como no livro \textit{Bagagem}.
\item O recurso autobiográfico é o mesmo recurso utilizado por Carrascoza, no livro \textit{Aos 7 e aos 40}.
\item A visão pessimista do amor opõe-se à valorização desse, como nos poemas otimistas de Adélia Prado.
\item A linguagem formal é utilizada tanto em \textit{Velha Infância}, como em \textit{Aos 7 e aos 40}, de Carrascoza.\end{enumerate}